\bookchapter{Capítulo 27: Para os pais também usarem com facilidade — Guia completo de acessibilidade para idosos no celular}

\chapterintro{Sinceramente, ao escrever este capítulo, pensei nos meus pais. Toda vez que aparece um pop-up de “atualização do sistema” no celular da minha mãe, ela não ousa clicar, com medo de estragar o aparelho. A fonte do WeChat do meu pai é menor que uma formiga, mas ele não sabe como aumentá-la, e ficou dois anos apertando os olhos para ler. Isso não é “burrice” deles — é que o celular, por padrão, trata todo mundo como um jovem de vinte e poucos anos. Neste capítulo, vou desatar, com a linguagem mais direta possível, os problemas mais angustiantes que os mais velhos enfrentam ao usar o celular.}

\booksection{Zero. Antes de tudo, ajuste a mentalidade: não é “idoso que não sabe usar”, é “celular que não foi bem feito”}

Antes de fazer qualquer configuração, é preciso esclarecer um conceito:\textbf{as interfaces dos smartphones atuais, desde o projeto, pouco consideram os usuários idosos}. Pop-ups minúsculos que mal se enxergam, gestos que exigem toque preciso, ícones quase idênticos — isso não é culpa do usuário, é falha do produto. Portanto, ao ajudar um idoso a usar o celular, a primeira coisa não é ensinar, é ajudar a\textbf{ajustar o ambiente}. Com o ambiente ajustado, a maior parte dos problemas desaparece.

Três regras de ouro ao ajudar idosos com o celular: (1)\textbf{Não apresse}— você usa smartphones há dez anos, eles podem estar usando há apenas três meses. (2)\textbf{Use analogias}— “o Wi-Fi é como o encanamento de casa: só sai água se a torneira estiver aberta”. Metáforas assim são mil vezes mais compreensíveis do que termos técnicos. (3)\textbf{Anote}— explicar verbalmente uma vez pode ser esquecido. Se você anotar os passos-chave num papel e colar na capinha do celular, isso vale mais que qualquer coisa.

\booksection{I. Primeiro passo: ative o “Modo Simplificado” — o celular fica mais limpo}

Os celulares atuais basicamente vêm com um “Modo Simplificado” integrado (alguns chamam de “Modo de Cuidado”, “Modo Sênior”, “Modo Idoso”). Ao ativá-lo, o celular inteiro se transforma:

\begin{itemize}
  \item \textbf{Fonte e ícones ficam bem grandes}— nunca mais será preciso usar óculos de leitura para ver o celular.
  \item \textbf{Ícones da tela inicial são drasticamente reduzidos}— mantém apenas os apps mais usados, como Telefone, SMS, Câmera, Galeria e WeChat. Todo o resto é recolhido.
  \item \textbf{A operação fica mais simples}— várias funções chamativas (tela negativa, assistente inteligente, janelas flutuantes, etc.) são desativadas por padrão. A interface prima pela limpeza visual.
  \item \textbf{O som fica mais alto}— toque de chamada, volume de notificações e som das ligações são automaticamente ajustados ao máximo.
\end{itemize}

\booksubsection{Como ativar o Modo Simplificado em cada marca (é só seguir os passos):}

\begin{center}
\begin{tabularx}{\linewidth}{|X|X|X|}
\hline
\textbf{Marca} & \textbf{Nome do modo} & \textbf{Como ativar} \tabularnewline
\hline
iPhone & Uso assistivo (não há um modo simplificado separado, mas é possível ajustar item por item) & Ajustes > Tela e Brilho > Tamanho do Texto, arraste ao máximo. Depois em "Acessibilidade > Tela e Tamanho do Texto", ative "Texto em Negrito" e "Formatos de Botão" \tabularnewline
Huawei/Honor & Modo Simplificado & Config. > Sistema e atualizações > Modo Simplificado > Ativar \tabularnewline
Xiaomi/Redmi & Modo Minimalista & Config. > Mais configurações > Modo Minimalista > Ativar \tabularnewline
OPPO/OnePlus & Modo Simplificado & Config. > Recursos especiais > Modo Simplificado \tabularnewline
vivo/iQOO & Tela Inicial Simplificada & Config. > Tela inicial, tela de bloqueio e papel de parede > Tela Inicial Simplificada \tabularnewline
Samsung & Modo Simplificado & Config. > Tela > Modo Simplificado \tabularnewline
\hline
\end{tabularx}
\end{center}

Se você usa outra marca, na barra de busca no topo das “Configurações” do celular, procure diretamente por “simplificado” ou “idoso”. Geralmente você encontra. Se não encontrar, siga o esquema de “ajuste manual” abaixo, item por item. O efeito é o mesmo.

\booksection{II. Se não quiser ativar o Modo Simplificado, ajuste manualmente — “Fonte grande + Ícones grandes” em cinco passos}

Alguns idosos não querem usar o Modo Simplificado, achando a interface muito “antiquada”. Nesse caso, ajuste manualmente os cinco itens a seguir. O efeito é igualmente bom:

\begin{enumerate}
  \item \textbf{Aumente a fonte ao máximo}: Config. > Tela > Tamanho da fonte, arraste até a extremidade direita. Aproveite e também ajuste o "Tamanho da tela" (alguns chamam de "Zoom da tela") ao máximo. Assim, todos os elementos da interface aumentam juntos.
  \item \textbf{Aumente os ícones}: na tela inicial, faça o gesto de pinça com dois dedos (como se fosse reduzir uma foto) para entrar nas configurações da tela inicial > Layout da tela inicial (ou Tamanho dos ícones). Escolha o layout com menos quadrículas (por exemplo, 3×5). Os ícones naturalmente ficarão maiores.
  \item \textbf{Ative "Texto de alto contraste"}: torna as bordas do texto mais nítidas. Essa opção pode ser encontrada em "Acessibilidade".
  \item \textbf{Desative notificações desnecessárias}: Config. > Notificações. Desative todas as notificações de apps, exceto Telefone, SMS e WeChat. Assim, o celular não fica apitando a toda hora, evitando incômodo.
  \item \textbf{Coloque os contatos frequentes na tela inicial}: pressione e segure um espaço vazio na tela inicial > Adicionar widget > Contatos. Coloque a discagem rápida do cônjuge, filhos, médico e contato de emergência na tela inicial. Um toque e já liga, sem precisar navegar pela lista de contatos.
\end{enumerate}

\begin{tcbpractice}{Mãos à obra: ajude um idoso a "decorar" o celular}
\begin{enumerate}
  \item Ative o Modo Simplificado (conforme a tabela de marcas acima).
  \item Ative a exibição da porcentagem da bateria (Config. > Bateria > Porcentagem da bateria), para que saibam sempre quanto de carga resta.
  \item Abra a barra de notificações e remova ou oculte os atalhos rápidos que os idosos não usam, como Bluetooth, NFC e Modo Avião.
  \item Coloque quatro discagens rápidas na tela inicial: cônjuge, filhos, médico de família (se houver) e síndico/administração do condomínio.
\end{enumerate}

\end{tcbpractice}

\booksection{III. Prevenção de fraudes — por favor, veja este capítulo junto com os mais velhos}

Idosos são o principal alvo de fraudes por telecomunicações. Não é que eles tenham baixa vigilância; é que o discurso dos golpistas é cuidadosamente elaborado, mirando especificamente nas vulnerabilidades psicológicas dos idosos. A seguir, os tipos de golpe mais comuns que organizei e as formas de enfrentá-los. Sugiro que você leia diretamente para os mais velhos, ou imprima e cole ao lado do telefone.

\booksubsection{Seis golpes comuns (se receber, desligue imediatamente):}

\begin{itemize}
  \item \textbf{(1) "Seu neto se meteu em encrenca, envie dinheiro urgente"}— o golpista diz que "seu neto/netinha se envolveu em confusão, foi preso pela polícia, precisa de fiança". Lembre-se:\textbf{qualquer ligação pedindo transferência para "liberar" alguém é golpe}. Órgãos públicos e judiciais não pedem dinheiro por telefone.
  \item \textbf{(2) "Seu cartão de saúde foi suspenso"}— SMS dizendo que o cartão de saúde está irregular e pedindo para clicar em um link para "reativar". Lembre-se:\textbf{o órgão de saúde ou previdência não envia SMS pedindo para clicar em links}. Nesses casos, ligue diretamente para o 12333 (linha direta oficial da previdência social) para verificar.
  \item \textbf{(3) "Investimento de alto rendimento, lucro garantido sem risco"}— em grupos de WeChat ou SMS, recomendam vários "investimentos para aposentadoria" ou "projetos apoiados pelo governo", prometendo altos retornos. Lembre-se:\textbf{se prometem rendimento anual acima de 5\% e ainda "capital garantido", 99,9\% é golpe}.
  \item \textbf{(4) "Sou seu chefe/velho amigo, mudei de número"}— fingem ser conhecidos pedindo dinheiro emprestado. Lembre-se:\textbf{ligar para confirmar resolve tudo}. Quem diz que "não pode atender agora" tem 80\% de chance de ser golpista.
  \item \textbf{(5) "Ovos grátis / palestra sobre saúde"}— eventos presenciais gratuitos que induzem os idosos a ouvir "palestras de saúde" e depois vendem suplementos a preços exorbitantes. Lembre-se:\textbf{não existe almoço grátis. Por trás de eventos gratuitos, sempre há um negócio}.
  \item \textbf{(6) "Seu celular está com vírus, precisa de reparo remoto"}— o golpista faz você baixar um app de controle remoto (como AnyDesk, TeamViewer) e então controla remotamente seu celular para transferir dinheiro. Lembre-se:\textbf{nunca permita que ninguém controle remotamente seu celular!}Fabricantes de celular, operadoras de telefonia e bancos nunca ligarão pedindo para você instalar software de controle remoto.
\end{itemize}

\booksubsection{Três linhas de defesa. Bem estabelecidas, bloqueiam 90\% dos golpes:}

\begin{enumerate}
  \item \textbf{Primeira linha: bloqueio nativo do celular}— Android: ative o "Bloqueio de spam" (no app Discador ou no app Gestor do Telefone). iPhone: ative "Silenciar Chamadas Desconhecidas" (Ajustes > Telefone > Silenciar Desconhecidos).
  \item \textbf{Segunda linha: bloqueio da operadora}— envie SMS "KTFSR" para 10086 (China Mobile) para ativar gratuitamente o bloqueio de chamadas de spam de alta frequência. China Telecom e China Unicom também têm serviços semelhantes. Ligue para o atendimento da operadora e peça para ativarem para você.
  \item \textbf{Terceira linha: App Nacional Antifraude}— baixe o app "Central Nacional Antifraude" na loja de aplicativos, registre-se e ative o alerta de chamadas. Este app está conectado ao banco de dados da polícia. O número do golpista é identificado assim que a chamada entra.
\end{enumerate}

\begin{tcbwarn}{"Mantra antifraude" para os mais velhos (imprima e cole ao lado do telefone)}
\begin{itemize}
  \item \textbf{Chamada de desconhecido pedindo transferência Desligue!}
  \item \textbf{SMS com link para clicar Apague!}
  \item \textbf{Promessa de alto rendimento com capital garantido Bloqueie!}
  \item \textbf{Pedindo para instalar software de controle remoto Desligue na hora!}
  \item \textbf{Policial de verdade não resolve caso por telefone, muito menos pede dinheiro!}
  \item \textbf{Na dúvida, ligue primeiro para seus filhos e pergunte.}
\end{itemize}

\end{tcbwarn}

\booksection{IV. Assistência remota — você pode ajudar seus pais a consertar o celular mesmo estando longe}

Às vezes, o problema que os mais velhos encontram é tão difícil de descrever por telefone que eles não acham o botão correspondente. É aí que a “Assistência Remota” entra em cena — você pode ver diretamente a tela do celular deles no seu próprio celular e ajudá-los a operar. Cada marca tem uma solução diferente:

\begin{itemize}
  \item \textbf{Huawei/Honor}: vem com o app "Chat Connect", que tem a função "Compartilhamento de Tela". Instale o app em ambos os celulares. Inicie uma videochamada e toque em "Compartilhar Tela". Você pode desenhar círculos na tela deles para indicar "clique aqui".
  \item \textbf{Xiaomi/Redmi}: o app de Videochamada Xiaomi tem "Assistência Remota", com compartilhamento de tela + chamada de voz simultâneos.
  \item \textbf{OPPO/vivo}: têm "Espaço Familiar" ou "Guardião Remoto" integrados. Procure por "remoto" nas configurações.
  \item \textbf{Solução universal (funciona para todas as marcas)}: baixe o app "AnyDesk" ou "TeamViewer" (a versão gratuita é suficiente). Instale em ambos os celulares. Registre a mesma conta ou adicione como amigo. Peça aos pais que ditem o código de verificação exibido na tela deles. Você poderá ver e operar remotamente a tela deles.
  \item \textbf{Usuários de iPhone}: o FaceTime suporta "Compartilhar Tela" durante a chamada. Durante uma chamada FaceTime, toque no botão de compartilhamento de tela para que a outra pessoa veja sua tela. Mas, por limitação de segurança do iOS, não é possível operar remotamente — apenas "ensinar por comando".
\end{itemize}

\booksection{V. “Adaptação para idosos” do WeChat — basta dizer estas frases para os mais velhos}

O WeChat é o app mais usado pelos idosos, mas sua interface é muito densa e confusa. Aqui estão alguns ajustes de efeito imediato:

\begin{itemize}
  \item \textbf{Aumente a fonte do próprio WeChat}: WeChat > Eu > Configurações > Geral > Tamanho da Fonte, arraste ao máximo. Atenção: isso afeta apenas a fonte dentro do WeChat, não a fonte do sistema do celular.
  \item \textbf{Ative o "Modo de Cuidado"}: WeChat > Eu > Configurações > Modo de Cuidado > Ativar. Após ativar, fontes e botões ficam maiores, o contraste de cores é reforçado e é possível "Ouvir mensagens de texto" com um toque (toque na mensagem de texto para ouvi-la em voz alta).
  \item \textbf{Ajude a renomear os contatos importantes}: na lista de contatos, renomeie filhos e parentes com nomes fáceis de reconhecer — por exemplo, "Filho mais velho - Xiaoming", "Filha do meio - Xiaohong". Assim, ao abrir o WeChat, eles sabem quem é quem.
  \item \textbf{Ensine a usar "Entrada de voz" em vez de digitar}: na caixa de chat, ao lado do teclado, mantenha pressionada a barra de espaço (Android) ou toque no ícone do microfone (iPhone). Fale diretamente e o texto será convertido e enviado. A precisão do reconhecimento já está muito alta. Para idosos que não sabem digitar em pinyin, é uma comodidade "revolucionária".
  \item \textbf{Desative funções não usadas na "Página Descobrir"}: WeChat > Eu > Configurações > Geral > Gerenciar Página Descobrir. Desative tudo o que não for usado em "Momentos", "Canais", "Transmissão ao Vivo", "Compras" e "Jogos" para reduzir distrações.
\end{itemize}

\booksection{VI. E se o celular for perdido? Ajude os mais velhos a fazer um “Plano de Emergência”}

Quem se prepara, vence. Antes que o celular seja perdido, faça estas coisas primeiro:

\begin{enumerate}
  \item \textbf{Confirme que "Buscar meu celular" está ativado}: iPhone: Ajustes > Seu nome > Buscar > Buscar iPhone. Android: procure "Buscar Dispositivo" nas Configurações. Certifique-se de que esta opção está ligada.
  \item \textbf{Anote as senhas e contas importantes em um papel}: a senha do Apple ID / conta Huawei / conta Xiaomi. Guarde uma cópia trancada na gaveta de casa e entregue outra a um filho de confiança. Depois de perder o celular, você precisará fazer login nessas contas para bloquear e localizar remotamente. Se esquecer a senha, será um problema.
  \item \textbf{Ensine uma frase aos mais velhos}:\textbf{"Perdeu o celular? A primeira coisa é ligar para a família."}A família pode acessar imediatamente o iCloud ou a página web Buscar Dispositivo do Android e marcar remotamente como "Modo Perdido". Assim, quem encontrar o celular não conseguirá abri-lo nem usá-lo.
\end{enumerate}

\booksection{VII. Compatibilidade com fones de ouvido sem fio e aparelhos auditivos — uma boa função frequentemente ignorada}

Muitos idosos têm problemas auditivos, mas não sabem que o celular pode funcionar em conjunto com aparelhos auditivos. O iPhone e alguns celulares Android suportam aparelhos auditivos MFi (Made for iPhone), podendo transmitir o som da chamada diretamente para o aparelho auditivo. Muito mais claro do que segurar o celular no ouvido.

\begin{itemize}
  \item \textbf{Como conectar}: iPhone: Ajustes > Acessibilidade > Dispositivos Auditivos. Ative o Bluetooth e siga as instruções de pareamento do aparelho auditivo.
  \item \textbf{Se não usar aparelho auditivo}: em "Acessibilidade > Áudio e Visual", ative "Cancelamento de Ruído do Telefone" e "Áudio Mono". Isso melhora a clareza das chamadas.
  \item \textbf{Android}: em "Acessibilidade > Melhoria Auditiva", há configurações semelhantes. A localização varia por marca. Procure diretamente por "auditivo" ou "audição".
\end{itemize}

\booksection{VIII. Funções relacionadas à saúde — o celular não é só uma ferramenta de comunicação}

\begin{itemize}
  \item \textbf{SOS de Emergência}: pressione o botão liga/desliga cinco vezes seguidas (na maioria dos celulares). O celular discará automaticamente para o número de emergência e enviará sua localização para os contatos de emergência predefinidos. Configure com antecedência para os mais velhos — iPhone: "Ajustes > SOS de Emergência". Android: "Config. > Segurança > Emergência". Antes de ensinar, faça um ensaio para que não fiquem nervosos e errem o botão na hora.
  \item \textbf{Detecção de queda}: o Apple Watch tem essa função. Detecta quando você cai e fica imóvel por um tempo, discando automaticamente para emergência e notificando a família. Se tiver orçamento, vale a pena dar um de presente para os mais velhos. Pode salvar vidas em momentos críticos.
  \item \textbf{Lembrete de medicamentos}: o app "Saúde" do iPhone tem a função integrada de lembrete de medicamentos. Configure o nome, horário e dosagem do remédio. Na hora certa, um alerta aparece. No Android, use "Alarme" ou baixe um app de "Lembrete de Remédios".
\end{itemize}

\begin{tcbsummary}{Resumo do capítulo}
\begin{itemize}
  \item Primeiro ative o Modo Simplificado ou ajuste manualmente fonte/ícones para maior. Com o ambiente ajustado, metade dos problemas desaparece
  \item Três regras de ouro ao ajudar idosos: não apressar, usar analogias, anotar
  \item Reforce repetidamente os seis golpes comuns. Três linhas de defesa bem estabelecidas bloqueiam 90\% dos golpes
  \item Assistência remota permite que você ajude seus pais a consertar o celular mesmo estando longe
  \item "Modo de Cuidado" + "Entrada de voz" são as duas funções mais práticas do WeChat para idosos
  \item SOS de Emergência e Detecção de Queda são configurações que salvam vidas. Vale a pena gastar cinco minutos para configurá-las
\end{itemize}

\end{tcbsummary}
